\documentclass[11pt]{article}
\usepackage[margin=1in]{geometry}
\usepackage{xcolor}
\usepackage{tcolorbox}
\usepackage{hyperref}
\usepackage{enumitem}
\usepackage{booktabs}
\usepackage{graphicx}

% Define colors
\definecolor{osaorange}{RGB}{220,115,38}
\definecolor{aiblue}{RGB}{30,144,255}

% Configure hyperref
\hypersetup{
    colorlinks=true,
    linkcolor=blue,
    urlcolor=blue,
    citecolor=blue
}

\title{\textbf{H524 Introduction to Biostatistics\\
\Large AI-Enhanced Group Data Analysis Project\\
Fall 2025}}
\author{Dr. John Molitor}
\date{}

\begin{document}
\maketitle

\begin{tcolorbox}[colback=aiblue!10,colframe=aiblue,boxrule=2pt,arc=2pt,left=6pt,right=6pt,top=6pt,bottom=6pt]
\centering
\textbf{\large Final Group Project Overview}\\[0.3em]
Teams of 3-4 students will conduct a statistical analysis of an assigned medical dataset, with a \textbf{focus on comparing AI tool outputs}. This project replaces the traditional final exam and emphasizes critical evaluation of AI-assisted statistical analysis (20\% of final grade).
\end{tcolorbox}

\section{Project Timeline}

\begin{itemize}[leftmargin=20pt]
\item \textbf{Week 3:} Team formation and dataset assignment
\item \textbf{Week 4:} Check-in - Data familiarization meeting
\item \textbf{Week 5-6:} AI tool comparison analysis phase
\item \textbf{Week 7:} Draft analysis due for peer review (optional)
\item \textbf{Week 8:} Final project submission
\item \textbf{Week 9:} Group presentations (15 minutes + 5 minutes Q\&A)
\end{itemize}

\section{Learning Objectives}

By completing this project, you will:
\begin{itemize}[leftmargin=20pt]
\item Apply descriptive and inferential statistical methods to real health data
\item Conduct appropriate hypothesis tests for your dataset
\item \textbf{Compare outputs from three different AI tools} (Claude, ChatGPT, Microsoft Copilot)
\item \textbf{Critically evaluate AI-generated statistical analyses}
\item Document and verify AI-generated results
\item Communicate statistical findings clearly and professionally
\end{itemize}

\section{Dataset Assignments}

Each team will be assigned ONE of the following datasets from the MedDataSets R package:

\subsection{Dataset 1: Anorexia Treatment Study}
\begin{itemize}
\item \textbf{File:} anorexia.csv
\item \textbf{Sample Size:} 72 patients
\item \textbf{Design:} Pre-post study with 3 treatment groups
\item \textbf{Variables:} Treatment group (Control, CBT, Family Therapy), Pre-weight, Post-weight
\item \textbf{Research Question:} Which treatment is most effective for weight gain in anorexia patients?
\item \textbf{Methods:} ANOVA, paired t-tests, change score analysis
\item \textbf{AI Divergence Expected:} High - Multiple valid analytical approaches
\end{itemize}

\subsection{Dataset 2: Tooth Growth Study}
\begin{itemize}
\item \textbf{File:} toothgrowth.csv
\item \textbf{Sample Size:} 60 guinea pigs
\item \textbf{Design:} 2 $\times$ 3 factorial (2 supplements $\times$ 3 doses)
\item \textbf{Variables:} Tooth length, Supplement type (OJ vs VC), Dose (0.5, 1.0, 2.0 mg)
\item \textbf{Research Question:} Does vitamin C promote tooth growth? Does delivery method matter?
\item \textbf{Methods:} One-way or two-way ANOVA, pairwise comparisons
\item \textbf{AI Divergence Expected:} High - Factorial design, interaction effects
\end{itemize}

\subsection{Dataset 3: Low Birth Weight Study}
\begin{itemize}
\item \textbf{File:} birthwt.csv
\item \textbf{Sample Size:} 189 births
\item \textbf{Design:} Observational case-control study
\item \textbf{Variables:} Low birth weight (binary), mother's age, weight, race, smoking, hypertension
\item \textbf{Research Question:} What factors are associated with low birth weight?
\item \textbf{Methods:} Chi-square tests, t-tests, risk ratios
\item \textbf{AI Divergence Expected:} High - Binary outcome, multiple predictors
\end{itemize}

\subsection{Dataset 4: Pima Indian Diabetes Study}
\begin{itemize}
\item \textbf{File:} pima\_diabetes.csv
\item \textbf{Sample Size:} 200 women
\item \textbf{Design:} Observational cross-sectional study
\item \textbf{Variables:} Diabetes status, pregnancies, glucose, BP, BMI, age
\item \textbf{Research Question:} What factors predict diabetes in this population?
\item \textbf{Methods:} T-tests, chi-square, comparison of means
\item \textbf{AI Divergence Expected:} High - Data quality issues (missing data coded as 0!)
\item \textbf{Special Note:} Tests AI data cleaning abilities
\end{itemize}

\subsection{Dataset 5: Type 2 Diabetes Treatment Trial}
\begin{itemize}
\item \textbf{File:} diabetes2.csv
\item \textbf{Sample Size:} 699 adolescents
\item \textbf{Design:} Randomized controlled trial
\item \textbf{Variables:} Treatment (lifestyle, metformin, metformin+), Outcome (success/failure)
\item \textbf{Research Question:} Which treatment is most effective for adolescent type 2 diabetes?
\item \textbf{Methods:} Chi-square test, risk ratios, multiple comparisons
\item \textbf{AI Divergence Expected:} Moderate - Cleaner dataset, but 3-group comparison
\end{itemize}

\section{Required Analysis Components}

\subsection{Part 1: Descriptive Analysis (20 points)}
\begin{itemize}
\item Data exploration and summary statistics
\item Appropriate visualizations (3-4 high-quality plots)
\item Check for missing data, outliers, and data quality issues
\item Assessment of data distributions and assumptions
\end{itemize}

\subsection{Part 2: Inferential Analysis (25 points)}
\begin{itemize}
\item Formulate clear research question and hypotheses
\item Select appropriate statistical test(s) based on data type and design
\item Check and report assumption testing (normality, homogeneity of variance)
\item Conduct hypothesis tests with proper interpretation
\item Report effect sizes and confidence intervals
\item Discuss practical vs. statistical significance
\end{itemize}

\subsection{Part 3: AI Tool Comparison (30 points) --- \textcolor{aiblue}{CRITICAL COMPONENT}}

\textbf{This is the core learning objective of the project!}

For your assigned dataset, you must:

\subsubsection{Step 1: Consult Three AI Tools}

Upload your dataset and ask for analytical guidance from:
\begin{enumerate}
\item \textbf{Claude} (claude.ai) - Upload CSV, ask for analysis
\item \textbf{ChatGPT} (chat.openai.com) - Upload CSV, ask for analysis
\item \textbf{Microsoft Copilot} (In RStudio) - Use for R code generation
\end{enumerate}

\textbf{Suggested prompt:}
\begin{verbatim}
"I have a dataset on [brief description]. Please suggest
appropriate statistical methods to analyze [research question].
Check assumptions and provide R code for the analysis."
\end{verbatim}

\subsubsection{Step 2: Document AI Responses}

Create a comparison table:

\begin{center}
\begin{tabular}{|p{2.5cm}|p{3.5cm}|p{3.5cm}|p{3.5cm}|}
\hline
\textbf{Aspect} & \textbf{Claude} & \textbf{ChatGPT} & \textbf{Copilot} \\
\hline
Method Suggested & & & \\
\hline
Assumptions Checked & & & \\
\hline
Effect Sizes Reported & & & \\
\hline
Code Quality & & & \\
\hline
Interpretation Style & & & \\
\hline
\end{tabular}
\end{center}

\subsubsection{Step 3: Critical Evaluation}

Answer these questions in your report:
\begin{enumerate}
\item \textbf{Agreement:} Did all three AI tools suggest the same statistical method? If not, why did they differ?
\item \textbf{Rigor:} Which AI tool was most thorough in checking assumptions?
\item \textbf{Errors:} Did any AI tool make mistakes? What were they?
\item \textbf{Verification:} How did you verify the AI outputs were correct?
\item \textbf{Best Approach:} Which AI tool's approach did you ultimately follow? Why?
\item \textbf{Synthesis:} Did you create a "best of all three" hybrid approach?
\end{enumerate}

\subsubsection{Step 4: Verification Process}

Document how you verified AI-generated results:
\begin{itemize}
\item Checked against course textbook and lecture notes
\item Manually verified key calculations
\item Ran AI-suggested code and examined output
\item Consulted office hours or discussion board
\item Cross-referenced with multiple AI tools
\end{itemize}

\subsection{Part 4: Final Report and Presentation (25 points)}

\textbf{Written Report (5-6 pages):}
\begin{itemize}
\item Introduction with research question (0.5 page)
\item Methods: Data description and analytical approach (1 page)
\item Results: Tables, figures, and statistical findings (2 pages)
\item \textbf{AI Comparison Section} (1.5 pages) --- \textcolor{aiblue}{Required!}
\item Discussion: Interpretation, limitations, conclusions (1 page)
\item Appendix: AI usage log and R code
\end{itemize}

\textbf{Group Presentation (15 minutes):}
\begin{itemize}
\item Overview of dataset and research question (2 min)
\item Key statistical findings (5 min)
\item \textbf{AI tool comparison highlights} (5 min)
\item Lessons learned about AI-assisted analysis (3 min)
\item Q\&A (5 min)
\end{itemize}

\section{Deliverables}

\subsection{Check-In (Week 4) - See separate document}
\begin{itemize}
\item Team meeting confirmation
\item Dataset successfully loaded in R
\item Preliminary descriptive statistics
\item Initial data exploration
\end{itemize}

\subsection{Final Submission (Week 8)}

Submit via Canvas:

\begin{enumerate}
\item \textbf{Written Report (.pdf):} 5-6 pages following structure above
\item \textbf{R Script (.R):} Well-commented, reproducible code
\item \textbf{AI Comparison Document (.pdf):}
   \begin{itemize}
   \item Screenshots of AI tool interactions
   \item Comparison table
   \item Critical evaluation
   \end{itemize}
\item \textbf{Presentation Slides (.pdf):} 12-15 slides
\item \textbf{Dataset (.csv):} Your cleaned version (if modifications made)
\item \textbf{Team Contribution Statement:} Who did what
\end{enumerate}

\section{Grading Rubric}

\begin{center}
\begin{tabular}{lcc}
\toprule
\textbf{Component} & \textbf{Weight} & \textbf{Points} \\
\midrule
Descriptive Analysis & 20\% & 20 \\
Inferential Analysis & 25\% & 25 \\
AI Tool Comparison & 30\% & 30 \\
Report \& Presentation & 25\% & 25 \\
\midrule
\textbf{Total} & \textbf{100\%} & \textbf{100} \\
\bottomrule
\end{tabular}
\end{center}

\subsection{Detailed Grading Criteria}

\subsubsection{Excellent (90-100\%)}
\begin{itemize}
\item Comprehensive and accurate statistical analysis
\item \textbf{Thorough AI comparison with specific examples of differences}
\item \textbf{Clear documentation of verification process}
\item Professional presentation with compelling visualizations
\item All team members contributed meaningfully
\end{itemize}

\subsubsection{Good (80-89\%)}
\begin{itemize}
\item Mostly correct statistical analysis
\item \textbf{Good AI comparison covering main differences}
\item Clear presentation with good visualizations
\item Evidence of verification, though not exhaustive
\end{itemize}

\subsubsection{Satisfactory (70-79\%)}
\begin{itemize}
\item Basic statistical analysis with minor errors
\item \textbf{AI comparison present but superficial}
\item Adequate presentation meeting requirements
\item Limited evidence of critical evaluation
\end{itemize}

\subsubsection{Needs Improvement (<70\%)}
\begin{itemize}
\item Major statistical errors or inappropriate methods
\item \textbf{Minimal or no AI comparison}
\item Poor presentation or missing components
\item Little evidence of verification or critical thinking
\end{itemize}

\section{Team Formation and Roles}

Each team of 3-4 students should designate:
\begin{itemize}
\item \textbf{Data Manager:} Lead data cleaning and initial exploration
\item \textbf{Statistical Lead:} Guide analytical approach and interpretation
\item \textbf{AI Comparison Lead:} Coordinate AI tool consultations and comparison
\item \textbf{Presentation Lead:} Coordinate slides and delivery
\end{itemize}

\textbf{Note:} All team members must contribute to all aspects. Roles are for coordination only.

\section{Resources and Support}

\subsection{Technical Resources}
\begin{itemize}
\item R and RStudio (required)
\item AI Tools: Claude (claude.ai), ChatGPT (chat.openai.com), Microsoft Copilot
\item MedDataSets R package (datasets also provided as CSV)
\item Course textbook: Pagano \& Gauvreau, \textit{Principles of Biostatistics}
\end{itemize}

\subsection{Getting Help}
\begin{itemize}
\item Office Hours: Wednesday 10:00-11:30 AM
\item Canvas Discussion Board for project questions
\item Special project support sessions: Weeks 5, 6, 7
\item Email for team meeting scheduling
\end{itemize}

\subsection{Recommended R Packages}
\begin{verbatim}
library(tidyverse)    # Data manipulation
library(ggplot2)      # Visualization
library(stats)        # Statistical tests
library(car)          # Assumption checks
library(effectsize)   # Effect size calculations
\end{verbatim}

\section{Academic Integrity}

\begin{tcolorbox}[colback=red!5,colframe=red,boxrule=1pt,arc=2pt,left=6pt,right=6pt,top=6pt,bottom=6pt]
\textbf{Important Reminders:}
\begin{itemize}
\item AI-generated content must be verified and documented
\item Each team member must contribute substantively
\item Plagiarism will result in project failure
\item Teams may discuss approaches, but analyses must be independent
\item \textbf{You are responsible for correctness of all AI-generated content}
\end{itemize}
\end{tcolorbox}

\section{Tips for Success}

\begin{enumerate}
\item \textbf{Start Early:} Don't wait until Week 7 to consult AI tools
\item \textbf{Meet Regularly:} Schedule weekly team meetings
\item \textbf{Document Everything:} Keep detailed notes of AI interactions
\item \textbf{Verify, Verify, Verify:} Never trust AI output without checking
\item \textbf{Ask Questions:} Use office hours if AI tools disagree
\item \textbf{Focus on Comparison:} The AI comparison is 30\% of your grade!
\item \textbf{Practice Presentation:} Rehearse timing and transitions
\end{enumerate}

\section{What Makes This Project Different}

\begin{tcolorbox}[colback=aiblue!5,colframe=aiblue,boxrule=1pt,arc=2pt,left=6pt,right=6pt,top=6pt,bottom=6pt]
\textbf{Traditional Biostatistics Project:}
\begin{itemize}
\item Analyze data
\item Report results
\item Done
\end{itemize}

\textbf{This AI-Enhanced Project:}
\begin{itemize}
\item Analyze data \textit{using three different AI tools}
\item Compare and contrast AI suggestions
\item Critically evaluate which approach is best
\item Verify and synthesize results
\item \textbf{Learn to work WITH AI, not just use it}
\end{itemize}
\end{tcolorbox}

This project prepares you for the reality of modern statistical practice, where AI tools are powerful assistants but require critical human oversight.

\end{document}
